\documentclass[11pt,a4paper]{article}

% ==========================
%      PACKAGES DE BASE
% ==========================
\usepackage[utf8]{inputenc}
\usepackage[T1]{fontenc}
\usepackage[french]{babel}
\usepackage{lmodern}
\usepackage{geometry}
\usepackage{amsmath}
\usepackage{amssymb}
% \usepackage{amsfonts}
\geometry{margin=2.5cm}
\usepackage{setspace}
\onehalfspacing

\usepackage{xcolor}
\usepackage{hyperref}
\hypersetup{
    colorlinks=true,
    linkcolor=blue!50!black,
    urlcolor=blue!70!black,
    citecolor=blue!50!black
}

\usepackage{lmodern}

% ==========================
%        TCOLORBOX
% ==========================
\usepackage{tcolorbox}
\tcbuselibrary{skins,breakable}

% Style global des boîtes
\tcbset{
    boxrule=0.8pt,
    arc=3pt,
    colback=white,
    colframe=black,
}

% Boîte Question
\newtcolorbox{questionbox}{
    colback=blue!3,
    colframe=blue!60!black,
    title=Question,
    fonttitle=\bfseries,
    breakable,
    left=6pt,right=6pt,top=6pt,bottom=6pt
}

% Boîte Réponse
\newtcolorbox{answerbox}{
    colback=green!3,
    colframe=green!50!black,
    title=Réponse,
    fonttitle=\bfseries,
    breakable,
    left=6pt,right=6pt,top=6pt,bottom=6pt
}

% Environnement QA : \begin{QA}{Texte de la question} ... \end{QA}
\newenvironment{QA}[1]{
    \begin{questionbox}
        \textbf{#1}
    \end{questionbox}
    \begin{answerbox}
}{
    \end{answerbox}
}

% ==========================
%    TITRE / AUTEUR / DATE
% ==========================
\title{\textbf{Préparation aux entretiens Quant}\\[0.5em]
\large Banque d'Investissement \,•\, Hedge Fund \,•\, Asset Management}
\author{Ton Nom}
\date{\today}

\begin{document}

% ==========================
%        PAGE DE GARDE
% ==========================
\begin{titlepage}
    \centering
    \vspace*{3cm}
    
    {\Huge \textbf{Préparation aux entretiens Quant}\par}
    \vspace{1cm}
    {\large Banque d'Investissement • Hedge Fund • Asset Management\par}
    
    \vspace{3cm}
    {\large Dossier de révision personnel}\par
    \vspace{0.5cm}
    {\large \textbf{Auteur :} Mohamed Lemine}\par
    \vspace{0.3cm}
    {\large \textbf{Dernière mise à jour :} \today}\par
    
    \vfill
    
    {\small
    Ce document regroupe des questions de cours, de probabilités, de modèles stochastiques,\\
    de produits dérivés et de brainteasers mathématiques, avec réponses détaillées.\\[0.5em]
    Utilisation : révision active avant les entretiens Quant.
    }
    
\end{titlepage}

% ==========================
%         SOMMAIRE
% ==========================
\tableofcontents
\newpage



\section{Monté Carlo}

\subsection{Question 1 :}


\[
Y_{\inf\{ i \in \mathbb{N} : U_i < \tfrac{f_X(Y_i)}{k\, f_Y(Y_i)} \}}
\sim X,
\qquad
\text{où } Y_i \sim Y,\ U_i \sim \mathcal{U}(0,1),\ k f_Y(x) \ge f_X(x).
\]



% ==========================
%      CONTENU EXEMPLE
% ==========================

\section{Question 01 :}

% \subsubsection*{Question :} Définir le mouvement brownien standard ? 
\begin{questionbox}
Définir le mouvement brownien standard ? 
\end{questionbox}


\subsubsection*{Reponse : }
Un mouvement brownien standard $(W_t)_{t \ge 0}$ est un processus stochastique vérifiant :
\begin{itemize}
    \item $W_0 = 0$ presque sûrement,
    \item accroissements indépendants,
    \item pour $t > s$, l'accroissement $W_t - W_s \sim \mathcal{N}(0, t-s)$,
    \item trajectoires continues presque sûrement.
\end{itemize}
C’est la brique de base de la modélisation continue en finance (Black–Scholes, EDS, formule d’Itô, etc.).
\newpage 

\section{Question 02 :}

\begin{questionbox}
Quelles sont les conditions pour qu'une variable aléatoire $Z$ soit une densité de changement de probabilité ?
    
\end{questionbox}

\subsubsection*{Reponse : }
Pour que $Z$ définisse un changement de probabilité de $\mathbb{P}$ vers une nouvelle mesure $\mathbb{Q}$ via :

\[
\frac{d\mathbb{Q}}{d\mathbb{P}} = Z,
\]

il faut que $Z$ vérifie les conditions suivantes :

\begin{enumerate}
    \item \textbf{$Z \ge 0$ presque sûrement}  
    $Z$ doit être une densité, donc ne peut pas être négative.

    \item \textbf{$\mathbb{E}_{\mathbb{P}}[Z] = 1$}  
    Ceci garantit que $\mathbb{Q}$ est une probabilité (masse totale égale à 1).

    \item \textbf{$Z$ est $\mathcal{F}$-mesurable}  
    Pour que $\mathbb{Q}$ soit bien définie sur le même espace probabilisable.

\end{enumerate}

Ainsi, $Z$ est une densité de Radon–Nikodym si et seulement si :

\[
Z \ge 0, \qquad Z \in L^1(\mathbb{P}), \qquad \mathbb{E}_{\mathbb{P}}[Z] = 1.
\]

Ce sont les conditions classiques utilisées dans les changements de mesure (ex : Girsanov).
\newpage

\begin{QA}{Quelles sont les conditions pour qu'un processus $(Z_t)_{t\ge 0}$ définisse un changement de mesure ?}
Pour que $(Z_t)$ définisse un changement de probabilité de $\mathbb{P}$ vers une mesure
$\mathbb{Q}$ via :
\[
\left.\frac{d\mathbb{Q}}{d\mathbb{P}}\right|_{\mathcal{F}_t} = Z_t,
\]
il faut que $(Z_t)$ vérifie les conditions suivantes :

\begin{enumerate}
    \item \textbf{Positivité :}  
    \[
    Z_t \ge 0 \quad \text{p.s.}
    \]
    Une densité ne peut pas être négative.

    \item \textbf{Intégrabilité :}  
    \[
    Z_t \in L^1(\mathbb{P}) \quad \forall t.
    \]

    \item \textbf{Normalisation :}  
    \[
    \mathbb{E}_{\mathbb{P}}[Z_t] = 1 \quad \forall t.
    \]
    Cela garantit que $\mathbb{Q}$ est bien une mesure de probabilité.

    \item \textbf{Condition martingale (essentielle) :}  
    \[
    (Z_t) \text{ est une martingale positive sous } \mathbb{P}.
    \]
    Cela assure la cohérence temporelle du changement de mesure :
    \[
    Z_t = \mathbb{E}_{\mathbb{P}}[Z_T \mid \mathcal{F}_t].
    \]

    \item \textbf{(Optionnel) Strictement positif :}  
    Si l’on veut que $\mathbb{P}$ et $\mathbb{Q}$ soient équivalentes :
    \[
    Z_t > 0 \quad \text{p.s.}
    \]
\end{enumerate}

Ainsi, un processus $(Z_t)$ est une densité locale de Radon–Nikodym s'il est positif,
intégrable, normalisé, et surtout une martingale. C'est le cadre classique du théorème
de Girsanov.
\end{QA}



% \section{Processus stochastiques}

% \begin{QA}{Définir le mouvement brownien standard.}
% Un mouvement brownien standard $(W_t)_{t \ge 0}$ est un processus stochastique vérifiant :
% \begin{itemize}
%     \item $W_0 = 0$ presque sûrement,
%     \item accroissements indépendants,
%     \item pour $t > s$, l'accroissement $W_t - W_s \sim \mathcal{N}(0, t-s)$,
%     \item trajectoires continues presque sûrement.
% \end{itemize}
% C’est la brique de base de la modélisation continue en finance (Black–Scholes, EDS, formule d’Itô, etc.).
% \end{QA}

% \begin{QA}{Que signifie « accroissements indépendants et stationnaires » ?}
% \textbf{Accroissements indépendants} : pour tout $0 \le t_0 < t_1 < \dots < t_n$,
% les variables $W_{t_1}-W_{t_0},\, W_{t_2}-W_{t_1},\, \dots,\, W_{t_n}-W_{t_{n-1}}$ sont indépendantes.\\[0.5em]
% \textbf{Accroissements stationnaires} : la loi de $W_{t+h} - W_t$ ne dépend que de $h$ (la longueur de l’intervalle), pas de $t$.
% \end{QA}

% \section{Modèle de Black--Scholes}

% \begin{QA}{Donner la formule du prix d’un call européen dans le modèle de Black--Scholes.}
% Le prix d'un call européen de strike $K$ et maturité $T$ est :
% \[
% C_0 = S_0 N(d_1) - K e^{-rT} N(d_2),
% \]
% avec :
% \[
% d_1 = \frac{\ln(S_0/K) + (r + \tfrac{1}{2}\sigma^2)T}{\sigma\sqrt{T}},
% \qquad
% d_2 = d_1 - \sigma\sqrt{T},
% \]
% où :
% \begin{itemize}
%     \item $S_0$ : prix spot de l’actif sous-jacent,
%     \item $K$ : strike,
%     \item $r$ : taux sans risque,
%     \item $\sigma$ : volatilité,
%     \item $N(\cdot)$ : fonction de répartition de la loi normale standard.
% \end{itemize}
% \end{QA}

% \section{Brainteasers / Probabilités}

% \begin{QA}{On lance deux dés à six faces. Quelle est la probabilité que la somme soit égale à 7 ?}
% Les couples $(1,6),(2,5),(3,4),(4,3),(5,2),(6,1)$ donnent une somme de $7$, soit $6$ issues favorables.\\
% Il y a $6 \times 6 = 36$ issues possibles au total.\\[0.3cm]
% Donc
% \[
% \mathbb{P}(\text{Somme} = 7) = \frac{6}{36} = \frac{1}{6}.
% \]
% \end{QA}

% ==========================
%   FIN DU DOCUMENT
% ==========================

\end{document}
